\section*{Hoofdstuk \ref{ch:b1:matrices} --- Matrices}

\begin{solution}{exo:b1:matrices:1}
\[
AB = \begin{pmatrix} 2 & 1\\ 1 & 0\end{pmatrix},
\quad
BA = \begin{pmatrix} 0 & 1\\ 1 & 2\end{pmatrix},
\quad
A^2 - B^2 = \begin{pmatrix} 1 & 4\\ 0 & 1\end{pmatrix} - I
= \begin{pmatrix} 0 & 4\\ 0 & 0\end{pmatrix},
\]
\[
(A+B)(A-B) = A^2 - AB + BA - B^2
= \begin{pmatrix} 0 & 4\\ 0 & 0\end{pmatrix} +
\begin{pmatrix} -2 & 0\\ 0 & 2 \end{pmatrix}
= \begin{pmatrix} -2 & 4\\ 0 & 2\end{pmatrix}.
\]
Zij verschillen met $BA - AB \neq 0$: de identiteit $(a+b)(a-b) =
a^2 - b^2$ vereist commutativiteit, en die faalt hier.
\end{solution}

\begin{solution}{exo:b1:matrices:2}
$M$: $L_2 \leftarrow L_2 - 2L_1$ doodt de tweede rij; $L_3
\leftarrow L_3 - L_1$ geeft $(0, -1, -2)$. Twee spillen:
$\operatorname{rk} M = 2$.

$N$: $L_3 \leftarrow L_3 - L_1$ geeft $(0,1,1,1) = L_2$; daarna
$L_3 \leftarrow L_3 - L_2 = 0$. Twee spillen: $\operatorname{rk} N
= 2$.
\end{solution}

\begin{solution}{exo:b1:matrices:3}
Reductie van $(A \mid I_3)$: $L_2 \leftarrow L_2 - 2L_1$, $L_3
\leftarrow L_3 - L_1$:
\[
\begin{pmatrix}
1 & 0 & 1 & 1 & 0 & 0\\
0 & 1 & -1 & -2 & 1 & 0\\
0 & 1 & 0 & -1 & 0 & 1
\end{pmatrix}
\xrightarrow{L_3 \leftarrow L_3 - L_2}
\begin{pmatrix}
1 & 0 & 1 & 1 & 0 & 0\\
0 & 1 & -1 & -2 & 1 & 0\\
0 & 0 & 1 & 1 & -1 & 1
\end{pmatrix},
\]
daarna $L_1 \leftarrow L_1 - L_3$, $L_2 \leftarrow L_2 + L_3$:
\[
A^{-1} = \begin{pmatrix}
0 & 1 & -1\\
-1 & 0 & 1\\
1 & -1 & 1
\end{pmatrix}.
\]
\emph{Controle:} de eerste rij van $A$ maal de eerste kolom van
$A^{-1}$: $1 \cdot 0 + 0\cdot(-1) + 1\cdot 1 = 1$; maal de tweede
kolom: $1 - 0 - 1 = 0$; maal de derde: $-1 + 0 + 1 = 0$.
\end{solution}

\begin{solution}{exo:b1:matrices:4}
$u(1) = 1$, $u(X) = X + 1$, $u(X^2) = X^2 + 2X + 1$: de kolommen
\omterm{prop:b1:vspaces:coordinates}{coördinaten} in $(1, X, X^2)$ geven
\[
M = \begin{pmatrix}
1 & 1 & 1\\
0 & 1 & 2\\
0 & 0 & 1
\end{pmatrix}.
\]
$u$ is inverteerbaar omdat zij de voor de hand liggende inverse $P
\mapsto P(X - 1)$ heeft (samenstelling van substituties). Haar
matrix wordt op dezelfde wijze verkregen uit $u^{-1}(X^k) =
(X-1)^k$:
\[
M^{-1} = \begin{pmatrix}
1 & -1 & 1\\
0 & 1 & -2\\
0 & 0 & 1
\end{pmatrix}.
\]
\end{solution}

\begin{solution}{exo:b1:matrices:5}
$N = \begin{pmatrix} 0 & 1\\ 0 & 0\end{pmatrix}$, $N^2 = 0$. Omdat
$2I$ en $N$ commuteren, breekt de binomiale ontwikkeling na twee
termen af:
\[
A^k = (2I + N)^k = 2^k I + k\,2^{k-1} N
= \begin{pmatrix} 2^k & k\,2^{k-1}\\ 0 & 2^k\end{pmatrix}.
\]
(Controle bij $k = 2$: $A^2 = \begin{pmatrix}4 & 4\\ 0 &
4\end{pmatrix}$, juist volgens het rechtstreekse product.)
\end{solution}

\begin{solution}{exo:b1:matrices:6}
\omterm{def:b1:matrices:transpose}{Sporen}: $\operatorname{tr}(AB - BA) = \operatorname{tr}(AB) -
\operatorname{tr}(BA) = 0$ (\cref{def:b1:matrices:transpose}),
terwijl $\operatorname{tr}(I_n) = n \neq 0$ in $\R$ of $\C$. Geen
oplossing. (Op oneindigdimensionale ruimten \emph{is} de
identiteit wel realiseerbaar --- afleiden en vermenigvuldigen met
$x$ voldoen eraan --- precies omdat er daar geen \omterm{def:b1:matrices:transpose}{spoor} bestaat.)
\end{solution}

\begin{solution}{exo:b1:matrices:7}
Telescoperend product, waarbij alle machten van $A$ commuteren:
\[
(I - A)(I + A + \dots + A^{m-1}) = I - A^m = I ,
\]
en \cref{prop:b1:matrices:ring} tilt de eenzijdige inverse op. Voor
de toepassing: de gegeven matrix is $I + N$ met
\[
N = \begin{pmatrix} 0 & 2 & 3\\ 0 & 0 & 2\\ 0&0&0 \end{pmatrix},
\quad
N^2 = \begin{pmatrix} 0&0&4\\ 0&0&0\\ 0&0&0\end{pmatrix},
\quad N^3 = 0 ,
\]
dus, met $-N$ in plaats van $A$ in de formule:
\[
(I + N)^{-1} = I - N + N^2 =
\begin{pmatrix}
1 & -2 & 1\\
0 & 1 & -2\\
0 & 0 & 1
\end{pmatrix}.
\]
\end{solution}

\begin{solution}{exo:b1:matrices:8}
$A^2 = A$: het endomorfisme $a$ is een \omterm{def:b1:linmaps:projection}{projectie}
(\cref{thm:b1:linmaps:projchar}), $E = \operatorname{im} a \oplus
\ker a$ met $\dim\operatorname{im} a = r = \operatorname{rk} A$. In
een \omterm{def:b1:vspaces:free}{basis} die aan deze ontbinding is aangepast ($r$ vectoren van
het \omterm{def:b1:linmaps:kerim}{beeld}, daarna een \omterm{def:b1:vspaces:free}{basis} van de kern) is de matrix van $a$ gelijk
aan $\begin{pmatrix} I_r & 0\\ 0 & 0\end{pmatrix}$, met \omterm{def:b1:matrices:transpose}{spoor} $r$.
Het \omterm{def:b1:matrices:transpose}{spoor} is invariant onder \omterm{def:b1:matrices:changeofbasis}{basisverandering}:
$\operatorname{tr}(P^{-1}MP) = \operatorname{tr}(MPP^{-1}) =
\operatorname{tr} M$ volgens de cyclische identiteit. Bijgevolg is
$\operatorname{tr} A = r = \operatorname{rk} A$.
\end{solution}

\begin{solution}{exo:b1:matrices:9}
$J^2 = nJ$ (elke ingang van $J^2$ telt $n$ enen op). Zoek $M^{-1} =
\alpha I + \beta J$:
\[
(aI + bJ)(\alpha I + \beta J)
= a\alpha\, I + (a\beta + b\alpha + nb\beta)\, J .
\]
Dit is gelijk aan $I$ dan en slechts dan als $a\alpha = 1$ en
$a\beta + b\alpha + nb\beta = 0$, dat wil zeggen $\alpha =
\frac1a$ en $\beta(a + nb) = -\frac ba$. Is $a \neq 0$ en $a + nb
\neq 0$, dan
\[
M^{-1} = \frac 1a I - \frac{b}{a(a + nb)}\, J .
\]
Omgekeerd: is $a = 0$, dan heeft $M = bJ$ rang $\leq 1 < n$ (voor
$n \geq 2$): niet inverteerbaar ($n = 1$ is het scalaire geval). Is
$a + nb = 0$, dan voldoet de vector $v = (1, \dots, 1)^{\mathsf T}$
aan $Mv = (a + nb)v = 0$ met $v \neq 0$: niet inverteerbaar. Dus $M
\in GL_n \iff a \neq 0$ en $a + nb \neq 0$.
\end{solution}

\begin{solution}{exo:b1:matrices:10}
\emph{Som:} $\operatorname{im}(A + B) \subseteq \operatorname{im} A
+ \operatorname{im} B$ (elke $(A+B)x = Ax + Bx$), en Grassmann
begrenst de dimensie van een som door de som van de dimensies.

\emph{Sylvester:} zij $a$ de \omterm{def:b1:logic:map}{afbeelding} van $A$, beperkt tot $V =
\operatorname{im} B$ (van dimensie $\operatorname{rk} B$). Haar
\omterm{def:b1:linmaps:kerim}{beeld} is $\operatorname{im}(AB)$ (want $a(Bx) = ABx$), en de
dimensiestelling in $V$ geeft
\[
\operatorname{rk} B = \dim\ker(a_{|V}) + \operatorname{rk}(AB) .
\]
Nu is $\ker(a_{|V}) \subseteq \ker A$, van dimensie $n -
\operatorname{rk} A$: dus
\[
\operatorname{rk}(AB) \geq \operatorname{rk} B - (n -
\operatorname{rk} A) = \operatorname{rk} A + \operatorname{rk} B -
n . \qedhere
\]
\end{solution}

\begin{solution}{exo:b1:matrices:11}
\begin{enumerate}
  \item Per ingang is $(AD)_{ij} = a_{ij}\,d_j$ en $(DA)_{ij} =
        d_i\,a_{ij}$. Dus $AD = DA$ dan en slechts dan als
        $a_{ij}(d_j - d_i) = 0$ voor alle $i, j$; voor $i \neq j$ is
        de factor $d_j - d_i$ niet nul, wat $a_{ij} = 0$ afdwingt:
        $A$ is diagonaal. Omgekeerd commuteren diagonale matrices
        met elkaar.
  \item Commuteert $A$ met elke matrix, dan commuteert zij met
        $\operatorname{diag}(1, 2, \dots, n)$, dus is $A =
        \operatorname{diag}(\lambda_1, \dots, \lambda_n)$ volgens
        (1). Dan is $A E_{ij} = \lambda_i E_{ij}$ (alleen rij $i$
        van $E_{ij}$ overleeft) terwijl $E_{ij} A = \lambda_j
        E_{ij}$: commuteren met $E_{ij}$ dwingt $\lambda_i =
        \lambda_j$ af. Bijgevolg $A = \lambda I_n$; en scalaire
        matrices commuteren inderdaad met alles. Het centrum van
        $\mathcal{M}_n(K)$ is $K\,I_n$.
\end{enumerate}
\end{solution}

\begin{solution}{exo:b1:matrices:12}
\begin{enumerate}
  \item Is $\operatorname{rk} A = 1$, dan is het \omterm{def:b1:linmaps:kerim}{beeld} van $A$ een
        rechte $\operatorname{Vect}(C)$ met $C \neq 0$, dus is de
        $j$-de kolom van $A$ gelijk aan $\ell_j\,C$ voor scalairen
        $\ell_j$ (niet alle nul), dat wil zeggen $A = C L$ met $L =
        (\ell_1, \dots, \ell_n) \neq 0$. Omgekeerd, is $A = CL \neq
        0$, dan zijn alle kolommen veelvouden van $C$: rang $1$.
  \item $A^2 = C\,(L C)\,L$, en $LC$ is de scalair $\sum_i \ell_i
        c_i = \operatorname{tr}(CL) = \operatorname{tr} A$. Dus $A^2
        = (\operatorname{tr} A)\,A$, en met inductie $A^m =
        (\operatorname{tr} A)^{m-1} A$. Is $\operatorname{tr} A \neq
        0$, dan wordt geen enkele macht nul; is $\operatorname{tr} A
        = 0$, dan is $A^2 = 0$: een matrix van rang één is nilpotent
        dan en slechts dan als haar \omterm{def:b1:matrices:transpose}{spoor} nul is.
  \item Met $t = \operatorname{tr} A \neq -1$:
        \[
        (I_n + A)\Bigl(I_n - \frac{A}{1 + t}\Bigr)
        = I_n + A - \frac{A + A^2}{1 + t}
        = I_n + A - \frac{(1 + t)A}{1 + t} = I_n ,
        \]
        met $A^2 = tA$. Is $t = -1$, dan is $(I_n + A)A = A + A^2 =
        A - A = 0$ met $A \neq 0$, zodat $I_n + A$ elke kolom van
        $A$ ongelijk aan nul doodt: niet \omterm{def:b1:logic:inj}{injectief}, niet
        inverteerbaar.
\end{enumerate}
\end{solution}

\begin{solution}{pb:b1:matrices:1}
\textbf{1.} Euclidische deling van $X^n$ door de \omterm{def:b1:poly:def}{monische} $D$ van
graad $2$ (\cref{thm:b1:poly:division}): quotiënt en rest bestaan
en zijn eenduidig, en de rest heeft graad $\leq 1$: $X^n = Q_n D +
a_n X + b_n$. Voor $n = 0$: $Q_0 = 0$, $(a_0, b_0) = (0, 1)$; voor
$n = 1$: $(a_1, b_1) = (1, 0)$.

\textbf{2.} Vermenigvuldig met $X$ en reduceer $X^2 = D + sX - p$:
\[
X^{n+1} = X Q_n D + a_n X^2 + b_n X
= (X Q_n + a_n)\,D + (s\,a_n + b_n)\,X - p\,a_n .
\]
De laatste uitdrukking heeft de gedaante van een rest (graad $\leq
1$), dus geeft de eenduidigheid $a_{n+1} = s a_n + b_n$ en $b_{n+1}
= -p a_n$. Invullen van $b_{n+1} = -pa_n$ in $a_{n+2} = s a_{n+1} +
b_{n+1}$ geeft $a_{n+2} = s\,a_{n+1} - p\,a_n$.

\textbf{3.} Evalueer $X^n = Q_n D + a_n X + b_n$ in de wortels:
$\lambda^n = a_n\lambda + b_n$ en $\mu^n = a_n\mu + b_n$. Aftrekken
en delen door $\lambda - \mu \neq 0$ geeft
\[
a_n = \frac{\lambda^n - \mu^n}{\lambda - \mu},
\qquad
b_n = \lambda^n - a_n\lambda
= \frac{\lambda\mu^n - \mu\lambda^n}{\lambda - \mu} .
\]

\textbf{4.} In de dubbele wortel: $\lambda^n = a_n\lambda + b_n$.
Afleiden van de identiteit geeft $nX^{n-1} = Q_n'\,(X - \lambda)^2
+ 2Q_n\,(X - \lambda) + a_n$, en evaluatie in $\lambda$: $a_n =
n\lambda^{n-1}$; daarna $b_n = \lambda^n - n\lambda^{n} = (1 -
n)\lambda^{n}$.

\textbf{5.} Voor $P = \sum_i p_i X^i$ en $Q = \sum_j q_j X^j$ is
\[
P(M)\,Q(M) = \sum_{i,j} p_i q_j M^{i+j} = (PQ)(M),
\]
omdat de machten van de ene matrix $M$ met elkaar commuteren (de
sommen zijn duidelijk wegens de lineariteit). Is $D(M) = 0$, dan
geeft het invullen van $M$ in $X^n = Q_n D + a_n X + b_n$ dat $M^n
= Q_n(M)\,D(M) + a_n M + b_n I = a_n M + b_n I$.

\textbf{6.} Rechtstreekse producten:
\[
A^2 = \begin{pmatrix}
a^2 + bc & b(a + d)\\
c(a + d) & d^2 + bc
\end{pmatrix},
\qquad
s A = \begin{pmatrix}
a(a+d) & b(a+d)\\
c(a+d) & d(a+d)
\end{pmatrix},
\]
dus heeft $A^2 - sA$ nullen buiten de diagonaal en op de diagonaal
de waarden $a^2 + bc - a^2 - ad = bc - ad = -p$: $A^2 - sA + pI_2 =
0$.

\textbf{7.} Met $A' = \begin{pmatrix} a' & b'\\ c' &
d'\end{pmatrix}$ geeft het uitwerken van $p(AA') = (aa' + bc')(cb'
+ dd') - (ab' + bd')(ca' + dc')$: de termen $aa'cb'$ en $ab'ca'$
heffen elkaar op, de termen $bc'dd'$ en $bd'dc'$ eveneens, en wat
overblijft is
\[
aa'dd' - bca'd' + bcb'c' - adb'c'
= (ad - bc)(a'd' - b'c') = p(A)\,p(A').
\]
Is $p \neq 0$, dan geeft Cayley--Hamilton dat
$A\,\bigl(\tfrac1p(sI_2 - A)\bigr) = \tfrac1p(sA - A^2) = I_2$,
waaruit de inverse volgt (en \cref{prop:b1:matrices:ring} maakt
haar tweezijdig). Is $p = 0$ en zou $A$ inverteerbaar zijn, dan
geeft de multiplicativiteit $1 = p(I_2) = p(A)\,p(A^{-1}) = 0$:
onmogelijk. Dus $A \in GL_2 \iff p \neq 0$.

\textbf{8.} $s = 3$, $p = 2$, $D = X^2 - 3X + 2 = (X - 1)(X - 2)$:
$\lambda = 2$, $\mu = 1$, dus $a_n = 2^n - 1$ en $b_n = 2 - 2^n$
(vraag 3). Bijgevolg
\[
A^n = (2^n - 1)A + (2 - 2^n)I
= \begin{pmatrix} 1 & 2^n - 1\\ 0 & 2^n \end{pmatrix}.
\]
Controle: $A^2 = \begin{pmatrix} 1 & 3\\ 0 & 4\end{pmatrix}$,
zowel volgens de formule als door rechtstreeks te kwadrateren.

\textbf{9.} $s = 4$, $p = 3\cdot1 - 1\cdot(-1) = 4$: $D = X^2 - 4X
+ 4 = (X - 2)^2$, met dubbele wortel $\lambda = 2$. Vraag 4: $a_n =
n\,2^{n-1}$, $b_n = (1 - n)2^n$, dus
\[
A^n = n\,2^{n-1}A + (1 - n)2^n I
= 2^{n-1}\begin{pmatrix} n + 2 & n\\ -n & 2 - n
\end{pmatrix}.
\]
Bij $n = 2$: $2\begin{pmatrix} 4 & 2\\ -2 & 0\end{pmatrix} =
\begin{pmatrix} 8 & 4\\ -4 & 0 \end{pmatrix}$, en dat is $A^2$
rechtstreeks berekend.

\textbf{10.} Inductie: $F^1 = \begin{pmatrix} F_2 & F_1\\ F_1 &
F_0\end{pmatrix}$, en
\[
F^{n+1} = F^n F =
\begin{pmatrix} F_{n+1} + F_n & F_{n+1}\\
F_n + F_{n-1} & F_n \end{pmatrix}
= \begin{pmatrix} F_{n+2} & F_{n+1}\\ F_{n+1} & F_n
\end{pmatrix}.
\]
Hier is $s = 1$, $p = -1$, $D = X^2 - X - 1$ met wortels $\varphi$
en $\psi$ ($\varphi - \psi = \sqrt5$, $\varphi\psi = -1$). De rij
$(F_n)$ heeft $F_0 = 0 = a_0$, $F_1 = 1 = a_1$ en gehoorzaamt aan
dezelfde recurrentie als $(a_n)$: $F_n = a_n = (\varphi^n -
\psi^n)/\sqrt5$, de formule van Binet. Cassini: toepassing van de
multiplicativiteit uit vraag 7 op $F^n$ geeft
\[
F_{n+1}F_{n-1} - F_n^2 = p(F^n) = p(F)^n = (-1)^n .
\]

\textbf{11.} De voorwaarde is \omterm{def:b1:linmaps:def}{lineair} en bevat de nulrij: een
\omterm{def:b1:vspaces:subspace}{deelruimte}. Met inductie bepalen $u_0, u_1$ de rij $u$ \omterm{def:b1:linmaps:def}{lineair}, en
wordt elk paar beginwaarden door precies één oplossing
gerealiseerd: zoals in \cref{exo:b1:findim:10} wordt $E_D$
\omterm{def:b1:logic:inj}{bijectief} en \omterm{def:b1:linmaps:def}{lineair} geparametriseerd door $(u_0, u_1) \in K^2$:
$\dim E_D = 2$.

\textbf{12.} $(a_n)$ gehoorzaamt aan de recurrentie (vraag 2) met
$a_0 = 0$, $a_1 = 1$. Hetzelfde geldt voor $(b_n)$: $b_{n+2} =
-p\,a_{n+1} = -p(s a_n + b_n) = s\,b_{n+1} - p\,b_n$ (met twee keer
$b_{n+1} = -pa_n$), met $b_0 = 1$, $b_1 = 0$. De combinatie $v_n =
u_1 a_n + u_0 b_n$ is dan een oplossing met $v_0 = u_0$, $v_1 =
u_1$; twee oplossingen met dezelfde beginwaarden vallen samen
(inductie), dus is $u_n = u_1 a_n + u_0 b_n$ voor alle $n$.

\textbf{13.} $(\lambda^n)$ is een oplossing dan en slechts dan als
$\lambda^{n+2} = s\lambda^{n+1} - p\lambda^n$ voor alle $n$, dat
wil zeggen $D(\lambda) = 0$ (na deling door $\lambda^n \neq 0$;
merk op dat $\lambda, \mu \neq 0$ omdat $p = \lambda\mu \neq 0$).
Vrijheid van $\bigl((\lambda^n), (\mu^n)\bigr)$: een betrekking bij
$n = 0, 1$ geeft $c + c' = 0$ en $c\lambda + c'\mu = 0$, dus
$c(\lambda - \mu) = 0$: $c = c' = 0$. Twee \omterm{def:b1:vspaces:free}{vrije} vectoren in
dimensie $2$: een \omterm{def:b1:vspaces:free}{basis}. Dubbele wortel: $\bigl((n\lambda^n)\bigr)$
is een oplossing want, met $s = 2\lambda$ en $p = \lambda^2$,
\[
s(n+1)\lambda^{n+1} - p\,n\lambda^n
= \lambda^{n+2}\bigl(2(n+1) - n\bigr) = (n+2)\lambda^{n+2} ;
\]
de vrijheid bij $n = 0, 1$: $c = 0$, en daarna $c'\lambda = 0$ met
$\lambda \neq 0$.

\textbf{14.} $D = X^2 - X - 6 = (X - 3)(X + 2)$. Algemene
oplossing $u_n = A\,3^n + B(-2)^n$; de beginvoorwaarden geven $A +
B = 1$ en $3A - 2B = 8$, dus $A = 2$, $B = -1$:
\[
u_n = 2\cdot 3^n - (-2)^n .
\]
Controle: $u_2 = 18 - 4 = 14 = u_1 + 6u_0$; $u_3 = 54 + 8 = 62 =
u_2 + 6u_1 = 14 + 48$.

\textbf{15.} $C\begin{pmatrix} u_n\\ u_{n+1}\end{pmatrix} =
\begin{pmatrix} u_{n+1}\\ -p\,u_n + s\,u_{n+1}\end{pmatrix} =
\begin{pmatrix} u_{n+1}\\ u_{n+2}\end{pmatrix}$, en inductie geeft
de formule met $C^n$. Bovendien is $\operatorname{tr} C = 0 + s =
s$ en $p(C) = 0\cdot s - 1\cdot(-p) = p$: de begeleidende matrix
heeft precies $D$ als haar \omterm{def:b1:poly:def}{veelterm} van Cayley--Hamilton.

\textbf{16.} Voor elke oplossing $u$ van $u_{n+3} = \alpha u_{n+2}
+ \beta u_{n+1} + \gamma u_n$ voldoen de toestandsvectoren $v_n =
(u_n, u_{n+1}, u_{n+2})^{\mathsf T}$ aan $C_3 v_n = v_{n+1}$ (de
eerste twee rijen verschuiven, de laatste rij past de recurrentie
toe). Bijgevolg is
\[
D_3(C_3)\,v_0 = v_3 - \alpha v_2 - \beta v_1 - \gamma v_0 ,
\]
waarvan de drie componenten $u_{k+3} - \alpha u_{k+2} - \beta
u_{k+1} - \gamma u_k = 0$ zijn ($k = 0, 1, 2$). Omdat de
begintoestand $v_0 = (u_0, u_1, u_2)^{\mathsf T}$ \emph{heel} $K^3$
doorloopt (de beginwaarden zijn \omterm{def:b1:vspaces:free}{vrij}), doodt de matrix $D_3(C_3)$
elke vector: $D_3(C_3) = 0$.

\textbf{17.} Schrijf $X^n = Q\,D_3 + R_n$ met $\deg R_n \leq 2$ en
evalueer in elke wortel: $\lambda_i^n = R_n(\lambda_i)$. Dus is
$R_n$ een \omterm{def:b1:poly:def}{veelterm} van graad $\leq 2$ die de drie waarden
$\lambda_i^n$ in de drie verschillende knooppunten $\lambda_i$
interpoleert: wegens de eenduidigheid in
\cref{thm:b1:poly:lagrange} is $R_n = \sum_i \lambda_i^n L_i$ met
$(L_i)$ de \omterm{def:b1:vspaces:free}{Lagrange-basis} van de knooppunten. Invullen van $C_3$
(de vragen 5 en 16) geeft
\[
C_3^{\,n} = R_n(C_3) = \sum_{i=1}^{3} \lambda_i^n\,L_i(C_3),
\]
met de drie matrices $L_i(C_3)$ onafhankelijk van $n$: elke ingang
van $C_3^{\,n}$ is een vaste combinatie van $\lambda_1^n,
\lambda_2^n, \lambda_3^n$.

\textbf{18.} $D_3 = X^3 - 2X^2 - X + 2 = (X-1)(X+1)(X-2)$.
Algemene oplossing $u_n = A + B(-1)^n + C\,2^n$.
Beginvoorwaarden: $A + B + C = 0$, $A - B + 2C = 1$, $A + B + 4C =
1$. Aftrekken van de eerste van de derde: $3C = 1$, dus $C =
\frac13$; daarna $A + B = -\frac13$ en $A - B = \frac13$: $A = 0$,
$B = -\frac13$. Bijgevolg
\[
u_n = \frac{2^n - (-1)^n}{3}
\]
(de jacobsthalgetallen). Controle: $u_3 = \frac{8 + 1}{3} = 3 =
2u_2 + u_1 - 2u_0 = 2 + 1 - 0$.

\textbf{19.} Taylorontwikkeling van de \omterm{def:b1:poly:def}{veelterm} $X^n$ in
$\lambda$:
\[
X^n = \sum_{k=0}^{n} \binom nk \lambda^{n-k}(X - \lambda)^k ,
\]
en alle termen met $k \geq 3$ zijn \omterm{def:b1:arith:divides}{deelbaar} door $(X -
\lambda)^3$: de rest is
\[
R_n = \lambda^n + n\lambda^{n-1}(X - \lambda) + \binom
n2\lambda^{n-2}(X - \lambda)^2 .
\]
Voor $M = \lambda I + N$ met $N^3 = 0$ is $(M - \lambda I)^3 = N^3
= 0$, dus geeft vraag 5
\[
M^n = \lambda^n I + n\lambda^{n-1} N + \binom n2
\lambda^{n-2} N^2 ,
\]
en dat is precies de binomiale ontwikkeling van $(\lambda I +
N)^n$, afgebroken bij $N^2$ --- de twee methoden stemmen overeen.

\textbf{20.} De oplossingsruimte heeft dimensie $3$ (dezelfde
parametrisering door $(u_0, u_1, u_2)$ als in vraag 11), en elke
$(\lambda_i^n)$ is een oplossing. Vrijheid: stel
$c_1\lambda_1^n + c_2\lambda_2^n + c_3\lambda_3^n = 0$ voor $n = 0,
1, 2$. Leg $i$ vast en zij $L_i = \sum_{k \leq 2} p_k X^k$ de
\omterm{def:b1:poly:def}{veelterm} van Lagrange van de knooppunten met $L_i(\lambda_j) =
\delta_{ij}$. Dan is
\[
0 = \sum_{k=0}^{2} p_k\Bigl(\sum_j c_j\lambda_j^k\Bigr)
= \sum_j c_j\,L_i(\lambda_j) = c_i .
\]
Dus zijn alle $c_i = 0$: drie \omterm{def:b1:vspaces:free}{vrije} oplossingen in dimensie $3$,
een \omterm{def:b1:vspaces:free}{basis}; de algemene oplossing is $c_1\lambda_1^n +
c_2\lambda_2^n + c_3\lambda_3^n$.

\textbf{21.} Uit $F_k = F_{k+2} - F_{k+1}$ telescopeert de som:
\[
\sum_{k=1}^{n} F_k = \sum_{k=1}^{n}\bigl(F_{k+2} - F_{k+1}\bigr)
= F_{n+2} - F_2 = F_{n+2} - 1 .
\]

\textbf{22.} Neem de ingang $(1,2)$ van $F^{m+n} = F^m F^n$: het
linkerlid is $F_{m+n}$; het rechterlid is (rij $1$ van $F^m$) maal
(kolom $2$ van $F^n$), dat wil zeggen $F_{m+1}F_n + F_m F_{n-1}$.
Met $m = n$:
\[
F_{2n} = F_{n+1}F_n + F_nF_{n-1} = F_n\,(F_{n+1} + F_{n-1}).
\]

\textbf{23.} Volgens Binet is $F_n - \dfrac{\varphi^n}{\sqrt5} =
-\dfrac{\psi^n}{\sqrt5}$, en $\abs\psi = \frac{\sqrt5 - 1}2 < 1$,
dus
\[
\Bigl|F_n - \frac{\varphi^n}{\sqrt5}\Bigr|
\leq \frac{1}{\sqrt5} < \frac12
\qquad (n \geq 0):
\]
$F_n$ is het dichtstbijzijnde gehele getal bij
$\varphi^n/\sqrt5$.

\textbf{24.} $t_n = F_{n+1} + F_{n-1}$ is een combinatie van
verschoven fibonaccirijen en voldoet dus aan dezelfde recurrentie:
$t_{n+2} = t_{n+1} + t_n$; en $t_1 = F_2 + F_0 = 1$, $t_2 = F_3 +
F_1 = 3$: dit zijn de lucasgetallen $L_n$. De rij $\varphi^n +
\psi^n$ is een oplossing met dezelfde eerste twee waarden ($\varphi
+ \psi = 1$, $\varphi^2 + \psi^2 = (\varphi + \psi)^2 -
2\varphi\psi = 3$), dus $L_n = \varphi^n + \psi^n$. Ten slotte is
\[
F_n L_n = \frac{(\varphi^n - \psi^n)(\varphi^n +
\psi^n)}{\sqrt5} = \frac{\varphi^{2n} - \psi^{2n}}{\sqrt5} =
F_{2n},
\]
waarmee vraag 22 wordt teruggevonden.

\textbf{25.} (i) De vijf matrices $I, A, A^2, A^3, A^4$ leven in
het $4$-dimensionale $\mathcal{M}_2(K)$, dus doodt \emph{een}
\omterm{def:b1:poly:def}{veelterm} ongelijk aan nul van graad $\leq 4$ de matrix $A$; deel
II scherpte dit aan tot de expliciete kwadratische $A^2 = sA -
pI$, die alle machten in het vlak $\operatorname{Vect}(I, A)$
opsluit. (ii) De euclidische deling reduceert $X^n$ modulo die
kwadratische \omterm{def:b1:poly:def}{veelterm}, en de twee coëfficiënten van de rest
gehoorzamen aan de recurrentie met twee termen $a_{n+2} =
s\,a_{n+1} - p\,a_n$: machtsverheffen is iteratie geworden. (iii)
Vraag 6 is het geval $n = 2$ van de \emph{stelling van
Cayley--Hamilton}, die in elke dimensie geldt en in het volume van
bachelorjaar 2 wordt bewezen. (iv) De begeleidende matrix sluit de
kring: elke \omterm{pb:b1:matrices:1}{lineaire recurrentie} \emph{is} een macht van een
matrix, met dezelfde \omterm{def:b1:poly:def}{veelterm} $D$ die als gegevens van \omterm{def:b1:matrices:transpose}{spoor} en
determinant verschijnt, zodat het rekenen met resten recurrenties
oplost en machten berekent in één beweging.
\end{solution}
